\chapter{الأعداد العشرية}\label{ch:g5:decimals}

ظهرت الأعشار والأجزاء من مئة في \cref{ch:g4:decimals}؛ وإضافة
مستوى واحد آخر، وهو الأجزاء من ألف، تُتمّ النظام العشري. ويرسّخ
هذا الفصل قراءة الأعداد العشرية وتفكيكها --- وقبل كل شيء
مقارنتها --- دون الوقوع في فخاخها المشهورة.

\section{إلى الأجزاء من ألف}

\begin{definition}[المراتب العشرية]\label{def:g5:decimals:places}
بعد \omterm{def:g4:decimals:point}{النقطة العشرية} تأتي الأعشار، ثم الأجزاء من مئة، ثم
الأجزاء من ألف --- وكل منها يساوي عُشر ما قبله:
\[
4.362 = 4 + \frac{3}{10} + \frac{6}{100} + \frac{2}{1000}
= \frac{4\,362}{1000}.
\]
\end{definition}

\begin{omfigure}
\begin{tikzpicture}[scale=0.9]
  \draw (0,0) grid[xstep=2.3, ystep=0.85] (9.2,1.7);
  \node[font=\small] at (1.15,1.275) {آحاد};
  \node[font=\small] at (3.45,1.275) {أعشار};
  \node[font=\small] at (5.75,1.275) {أجزاء من مئة};
  \node[font=\small] at (8.05,1.275) {أجزاء من ألف};
  \node[font=\large, omDef] at (1.15,0.425) {$4$};
  \node[font=\large, omThm] at (3.45,0.425) {$3$};
  \node[font=\large, omThm] at (5.75,0.425) {$6$};
  \node[font=\large, omThm] at (8.05,0.425) {$2$};
  \fill (2.3,0.06) circle (2.5pt);
\end{tikzpicture}

{\small جدول المراتب للعدد $4.362$: الجزء الصحيح بالأزرق، والجزء
العشري بالأحمر، والنقطة بين الآحاد والأعشار.}
\end{omfigure}

\begin{example}[قراءات كثيرة، وعدد واحد]\label{ex:g5:decimals:readings}
يمكن قراءة $4.362$ «أربعة نقطة ثلاثة ستة اثنان»، أو «أربعة
و $362$ جزءًا من ألف»، أو تفكيكه $4 + 0.3 + 0.06 + 0.002$. وحشوه
بأصفار نهائية لا يغيّر شيئًا: $4.362 = 4.3620$. لكن
$4.362 \neq 4.0362$: \omterm{def:g1:counting:zero}{فالصفر} \emph{في الداخل} يدفع كل رقم إلى
مرتبة أصغر.
\end{example}

\begin{example}[أين يقع على المستقيم؟]\label{ex:g5:decimals:line}
لوضع $4.362$: بين $4$ و $5$؛ وبالتكبير، بين $4.3$ و
$4.4$؛ وبتكبير آخر، بين $4.36$ و $4.37$، أقرب إلى $4.36$.
فكل رقم عشري مستوى تكبير إضافي.
\end{example}

\begin{omfigure}
\begin{tikzpicture}[scale=1.15]
  \draw[-Stealth] (-0.2,1.6) -- (10.4,1.6);
  \foreach \i in {0,...,10} \draw (\i,1.52) -- (\i,1.68);
  \node[below, font=\small] at (0,1.5) {$4$};
  \node[below, font=\small] at (10,1.5) {$5$};
  \node[below, font=\small] at (3,1.5) {$4.3$};
  \node[below, font=\small] at (4,1.5) {$4.4$};
  \fill[omThm] (3.62,1.6) circle (1.7pt);
  \draw[omExo, thin] (3,1.45) -- (0,0.15);
  \draw[omExo, thin] (4,1.45) -- (10,0.15);
  \draw[-Stealth] (-0.2,0) -- (10.4,0);
  \foreach \i in {0,...,10} \draw (\i,-0.08) -- (\i,0.08);
  \node[below, font=\small] at (0,-0.1) {$4.3$};
  \node[below, font=\small] at (10,-0.1) {$4.4$};
  \node[below, font=\small] at (6,-0.1) {$4.36$};
  \fill[omThm] (6.2,0) circle (1.7pt)
    node[above, font=\small] {$4.362$};
\end{tikzpicture}

{\small التكبير على مستقيم الأعداد: المسافة من $4.3$ إلى $4.4$،
مكبَّرةً، مقطوعة هي نفسها إلى عشرة أجزاء من مئة --- ويقع $4.362$
بعد $4.36$ بقليل.}
\end{omfigure}

\section{مقارنة الأعداد العشرية}

\begin{method}[مقارنة عددين عشريين]\label{met:g5:decimals:compare}
\begin{enumerate}
  \item قارن الأجزاء الصحيحة: $6.1 > 5.987$.
  \item وعند تساويها: قارن الأجزاء العشرية رقمًا رقمًا
        ابتداءً من النقطة --- الأعشار، ثم الأجزاء من مئة، ثم
        الأجزاء من ألف؛
  \item والحشو بأصفار نهائية يجعل المقارنة عادلة:
        $2.5$ مقابل $2.48$ يصير $2.50$ مقابل $2.48$، و $50 > 48$
        جزءًا من مئة.
\end{enumerate}
والفخّ، للمرة الأخيرة: \emph{الجزء العشري الأطول لا يعني عددًا
أكبر} --- ففي $2.48$ أرقام أكثر مما في $2.5$ وهو أصغر منه.
\end{method}

\begin{example}\label{ex:g5:decimals:order}
رتّب $7.3$ و $7.09$ و $7.31$ و $7.299$. اِحشُ إلى ثلاثة أرقام
عشرية: $7.300$ و $7.090$ و $7.310$ و $7.299$. إذن
\[
7.09 < 7.299 < 7.3 < 7.31 .
\]
ولاحظ أن $7.299 < 7.3$ مع أن $299$ يبدو كبيرًا: فهو $299$ جزءًا
من ألف مقابل $300$ جزءًا من ألف.
\end{example}

\begin{example}[الحشر بين عددين عشريين]\label{ex:g5:decimals:between}
هل يوجد عدد بين $5.7$ و $5.8$؟ نعم، وكثير: $5.75$
و $5.71$ و $5.799$\,\dots\ وبين $5.79$ و $5.8$؟ وكثير أيضًا:
$5.795$ مثلًا. فبين عددين عشريين مختلفين يمكن \emph{دائمًا}
حشر عدد آخر --- ولا يوجد عدد عشري «تالٍ».
\end{example}

\section{تقريب الأعداد العشرية}

\begin{definition}[التقريب]\label{def:g5:decimals:round}
\omterm{def:g4:numbers:round}{للتقريب} إلى الوحدة (أو العُشر، أو الجزء من مئة)، اُنظر إلى الرقم
التالي: $5$ فما فوق يقرّب إلى الأعلى، وإلا فإلى الأسفل. فالعدد $4.362$
يقرَّب إلى $4$ (الوحدة)، وإلى $4.4$ (العُشر)، وإلى $4.36$ (الجزء من مئة).
\end{definition}

\begin{example}[النقود تقرَّب إلى القروش]\label{ex:g5:decimals:money}
الثمن المحسوب $7.4983$ يُعرض $7.50$: فالمبالغ في الحياة
الواقعية تقرَّب إلى الجزء من مئة. ولاحظ التتابع: $7.4983 \to
7.50$، حيث صار $49$ هو $50$ --- \omterm{def:g4:numbers:round}{فالتقريب} قد يتموّج عبر
عدة أرقام.
\end{example}

\section{تمارين}

\begin{exercise}[$\star$]\label{exo:g5:decimals:1}
اُكتب على صورة عدد عشري: $5 + \frac{2}{10} + \frac{7}{1000}$؛
\quad $\frac{85}{100}$؛ \quad $\frac{4\,507}{1000}$؛ \quad «اثنا عشر
وتسعة أجزاء من ألف».
\end{exercise}

\begin{exercise}[$\star$]\label{exo:g5:decimals:2}
في $27.418$: ماذا يعدّ الرقم $4$؟ والرقم $8$؟ وفكّك العدد
كما في \cref{def:g5:decimals:places}.
\end{exercise}

\begin{exercise}[$\star$]\label{exo:g5:decimals:3}
صواب أم خطأ؟ اِشرح بالمراتب.
$3.60 = 3.6$؛ \quad $0.5 = 0.05$؛ \quad $2.070 = 2.07$؛ \quad
$1.3 = 1.03$.
\end{exercise}

\begin{exercise}[$\star$]\label{exo:g5:decimals:4}
اُنقل وأكمل بالإشارة $<$ أو $>$ أو $=$:
\[
4.8 \;?\; 4.79, \qquad
0.302 \;?\; 0.32, \qquad
6.5 \;?\; 6.500, \qquad
9.09 \;?\; 9.1 .
\]
\end{exercise}

\begin{exercise}[$\star$]\label{exo:g5:decimals:5}
رتّب من الأصغر إلى الأكبر: $3.14$؛ \ $3.4$؛ \ $3.041$؛ \
$3.104$؛ \ $3.41$.
\end{exercise}

\begin{exercise}[$\star$]\label{exo:g5:decimals:6}
بين أي عددين صحيحين يقع كل عدد؟ وبين أي عددين برقم عشري
واحد؟ \quad $6.83$؛ \quad $0.492$؛ \quad
$12.06$.
\end{exercise}

\begin{exercise}[$\star$]\label{exo:g5:decimals:7}
قرِّب $8.276$: إلى الوحدة؛ وإلى العُشر؛ وإلى الجزء من مئة. وقرِّب
$3.95$ إلى العُشر.
\end{exercise}

\begin{exercise}[$\star$]\label{exo:g5:decimals:8}
أعطِ عددًا يقع تمامًا بين: $2.6$ و $2.7$؛ \quad $0.41$ و
$0.42$؛ \quad $5.99$ و $6$.
\end{exercise}

\begin{exercise}[$\star$]\label{exo:g5:decimals:9}
ارتفاعات أربع نباتات هي $0.85$ م و $1.2$ م و $0.9$ م و
$1.02$ م. رتّبها من الأقصر إلى الأطول.
\end{exercise}

\begin{exercise}[$\star\star$]\label{exo:g5:decimals:10}
تقول لُبنى: «$7.12 > 7.9$ لأن $12 > 9$». صحّح خطأها،
مقارنًا الأعشار أولًا.
\end{exercise}

\begin{exercise}[$\star\star$]\label{exo:g5:decimals:11}
جِد كل الأعداد ذات الرقمين العشريين بالضبط التي تقرَّب إلى
$6.4$ عند \omterm{def:g4:numbers:round}{التقريب} إلى العُشر. فما أصغرها وما
أكبرها؟
\end{exercise}
